\chapter{47}

\section{Bettmeralp/VS (CH), Montag, 27.
August}



«Du hast dich selbst übertroffen, Joh.»

Manuela blickte Joh anerkennend an, schloss die Augen und genoss Bissen
um Bissen des cremigen Safranrisottos. Auch Louis staunte. Joh kochte
wie ein Gott.

«Wo hast du kochen gelernt?»

Manuela schöpfte Louis eine grosse Portion frisch geernteten
Gartensalat.

«Ganz klar,» Joh wischte sich den Mund mit der Serviette ab, «von meiner
Grossmutter. Als Kind war ich mehr bei ihr als zuhause.»

Joh lachte dazu sonderbar, wie jemand, der aufgehört hat, zu hoffen.
Dennoch klang er weder zerknirscht noch leidend.

«Und mit acht, neun Jahren konnte ich mich selbständig verpflegen. Das
war wichtig, weil ich oft alleine war.»

«Ach.»

Louis konnte seinen Blick nicht von Joh abwenden. Diese Gelassenheit. Einmal mehr war ihm,
als käme Joh von woanders. Vielleicht hatte der sich über die Zeit 
einfach mit der Realität zu versöhnen gelernt oder sich ihr mindestens
wohlwollend angenähert.
Das Gespräch rund um Johs Kochgeschichte endete mit dem Abräumen des
Geschirrs. Louis hätte gerne mehr erfahren. War sich jedoch unsicher, ob
es angebracht war, weitere Fragen zu stellen. Er wusste nur zu gut, wie
lästig es sich anfühlen konnte, ausgefragt zu werden.

Zum Nachtisch setzten sie sich gemeinsam in die Sofaecke.
Louis zog ein kleines Notizbüchlein aus der Hosentasche und knipste
seinen Kugelschreiber an. Konzentriert und erwartungsvoll blickte er zu
Joh auf.

«Nun denn, lieber Ciro, erstmal zu deiner Alphütte.»

Joh klang dabei euphorisch. Er begann in einem Ordner zu blättern. Die
beiden ersten Seiten zeigten nahezu ein Dutzend Bilder. Louis rückte
näher zu Joh. Der erklärte ihm die Details. Von der Aussenansicht hätte
Louis niemals auf das geräumige Innere geschlossen. Die kleine Küche,
eine einfache Waschstation und die gemütliche Wohnecke befanden sich
unten. Ein schmaler Tisch mit vier zurechtgesägten Holzstämmen als
Stühle standen direkt neben einem schmalen Kochherd.
Im hinteren Bereich entdeckte Louis ein holzverschaltes Bett. Es sah
selbst geschreinert aus.

«Beim Betreten der Hütte stehst du direkt in der Küche, siehst du?»

Joh drehte das entsprechende Bild in Louis’ Richtung.
Ein schlichter Holzboden schien nachträglich in die Hütte gezogen worden
zu sein. Über eine steile, schmale Leiter gelangte man nach oben. Hier
befanden sich weitere Schlafplätze. Eine breite Matratze lag links und
eine schmalere rechts auf dem Boden. Helle Überwürfe lagen über Kissen und Decken. 
Der Raum war niedrig, stehen würde man nicht können. Aber zum Schlafen sah es einladend aus. Auf Bildern
des Unter- wie des Obergeschosses standen an verschiedenen Orten schwer
wirkende Glasvasen, in die grosse Sonnenblumen eingestellt waren.
Louis wunderte sich über die Einrichtung.

«Das ist eine grossartige zweite \emph{Schäferloft,»}. 
So einladend habe ich mir die Alphütte nicht vorgestellt.»

Joh hob die Augenbrauen und schielte zu Manuela.

«Die Hütte hat Manuela mit ihrem weiblichem Charme gestaltet. Bis wir
all das Material oben hatten, hat’s gedauert,» er verwarf die Hände,
«ein Freund hat die sperrigen, schweren Dinge mit dem Heli hochgeflogen.
Mit viel Glück konnte er den Flug mit einem anderen Einsatz irgendeinmal
verbinden.»

«Ja, es ist zwar gemütlich da oben, aber dennoch sehr einfach.»

Manuela war die Ergänzung offenbar wichtig.
Joh blätterte weiter. Das nächste Bild zeigte ein vergrössertes Foto
eines Strom-Generators, das die ganze Seite ausfüllte. Die Details der
Maschine waren gut zu erkennen.

«Wasser gibt es am Brunnen neben der Hütte, manchmal mehr, manchmal sehr
wenig. Bisher reichte es zum Kochen und zum Duschen aus. Bei dieser
Trockenheit musst du sehr sparsam damit umgehen,» fuhr Joh fort und
schaute ernst zu Louis hinüber.

«Duschen?»

Louis runzelte die Stirn.

«Ja, da, schau,» Joh suchte das entsprechende Bild fast am Ende der
Dokumentation, «gegen Süden gerichtet haben wir eine Solardusche
eingerichtet. Und klar, auch beim Duschen musst du sehr vorsichtig sein,
sonst stehst du nach dem Einseifen ohne Wasser da.»

Jetzt lachte Joh laut auf und Manuela lachte mit.

«Alles schon erlebt, Ciro.»

Sie grinste. Louis stellte sich die beiden dabei vor, wie sie mit einer
Seifenschicht überzogen nackt draussen standen und auf einen
Wasserstrahl warteten, der nicht kam.

«Wenn wir auch oben sind,» erklärte Joh weiter, «sind wir mehrheitlich
in dieser Hütte. Mazis Logis ist kleiner, hat auch Platz für drei
Personen, aber die Hütte ist zu dritt wirklich enorm eng. Sie befindet
sich ungefähr zwanzig Minuten zu Fuss von der grossen Hütte entfernt.
Weisst du, Ciro, die Schafe leben in kleinen Gruppen von höchstens zehn
Tieren. Und die Gruppen sind auf der Alp verteilt. Du und Mazi kümmert
euch je um drei davon. Deine Gruppen sind in der Nähe deiner Hütte am
Sömmern, die andern in Mazis Umgebung. Aber,» Joh machte eine kleine
Pause, «Mazi wird dich morgen genau in deine Arbeit einführen und dir
deine restlichen Fragen vor Ort beantworten, das ist einfacher, als wenn
ich dich nun mit weiteren Einzelheiten zutexe.»

Manuela war unterdessen aufgestanden und kam mit eisgekühltem Tee
zurück. Sie goss das Getränk in drei hohe Gläser und liess
Zitronenschnitze hineingleiten.
Während Joh Louis schliesslich einen wohl durchdachten Vortrag zum
Umgang mit dem Stromgenerator gab, sog Louis die kühle Erfrischung in
sich ein. Solarzellen erzeugten den Strom, der die Maschine antriebe, erklärte Joh.
Damit würden Louis’ Kühlbox und allfällige Kleingeräte wie sein Handy
betreiben können.

«Gekocht wird mit Gas. Die Flaschen bringt uns regelmässig ein Freund
mit dem Heli hoch.»

Der kleine Holzherd sei ein Relikt aus alter Zeit.

«Es ist nur wenig Holz oben und wir müssen jedes Scheit von hier unten
hochbringen. Die Waldgrenze ist bereits oberhalb unseres Hofes erreicht.
Aber das hast du bestimmt bemerkt. Darum, den Holzofen haben wir kaum je
in Betrieb.»

Louis nickte die Erklärungen nach und nach nur mehr ab. Er tat sich
schwer, konzentriert zu bleiben und fragte sich, ob Joh nicht zu viel
redete. Vieles verstand sich wirklich von selbst.
Bei dringendem Bedarf, erklärte Joh ungebremst weiter, könne Louis die
Energie auch für eine Stablampe nutzen. Wieder sprach er von gut
überlegter Nutzung. Die Erfahrung zeige, dass er damit auskommen könne.

«Und zusätzlich, Ciro, sind eine Petroleumlampe und Kerzen in der Hütte.
Auch die diversen Anschlusskabel sind oben. Denk an dein Handykabel.»

Mit einigen Stichworten hatte Louis notiert, was ihm wichtig erschien. 
Je mehr er wusste, desto klarer wurde ihm, seine Gitarre würde erstmal
unten bleiben. Louis gefiel die Aussicht, nach ein paar Tagen «Alp» für
zwei Nächte nach unten auf den Hof absteigen zu können.
Der Handy-Empfang, schloss Joh endlich, sei ganz gut. Joh legte den
Ordner auf den kleinen Tisch, stand auf und begab sich in die Küche.
«Kaffeezeit,» murmelte er leise. Er sah müde aus.
Louis streckte sich durch. Beim Anblick von Manuela ahnte er, dass eine Fortsetzung folgte. Sie setzte sich gerade hin.

«Wir haben ein System entwickelt, das wir vorwiegend an die Bedürfnisse
von Johs Jungs anpassten. Die meiste Zeit ist mindestens einer im
Programm auf dem Hof. Dass einer runter muss, wie gerade David,
geschieht selten. Manchmal betreuen wir zwei junge Männer gleichzeitig.
Ab Sonntag sind wir wieder im gewohnten Wochenwechsel,» Manuela blickte
kurz zu Joh hinüber, «sofern David wie geplant zurückkommt. Aber wie Joh
nach dem heutigen Gespräch meint, sieht es gut aus.»

Joh kam mit kleinen Espressotassen auf einem Servierbrett zurück.
Er reichte eine Louis und eine zweite Manuela und setzte sich Louis
gegenüber.
Manuela sprach weiter.  

«Am Sonntag kommen Joh und ich zusammen mit David rauf, bringen frische
Wäsche, Nahrungsmittel und lösen dich und Mazi für zwei Nächte ab. Du
und Mazi, ihr steigt ab und nehmt Abfall und Wäsche runter. Du bist dann
zusammen mit Mazi bis am Dienstag auf dem Hof. Ihr werdet den Abfall
entsorgen, die Wäsche waschen, Lieferungen entgegen nehmen und andere
Kleinigkeiten erledigen. Mazi weiss genau, was unten zu tun ist. Am
Dienstag früh steigt ihr wieder hoch und wenn ihr da seid, gehe ich
runter. Ich erwarte Besuch. Ihr seid dann die folgende Woche zu viert
oben.»

Manuela nippte an ihrem Espresso. Louis’ Kopf brummte.

«Aber ab morgen, Ciro, bis Sonntag, bist du ein Glückspilz. Während
dieser Tage gehört die Hütte dir,» lachte sie, «die Woche später, wie
gesagt, wohnen Joh, David und du da.»

Ihr Konzept leuchtete Louis ein. Definitiv, nachdem Joh ihm seinen
Auftrag als Integrationsbegleiter detailliert geschildert hatte. Wo die
Jungs waren, war in der Regel auch Joh, ob unten auf dem Hof oder oben
auf der Alp. Mazi hatte sich offenbar zum perfekten Mitarbeiter im Team
gemausert. Er schien grosses Vertrauen zu geniessen und unentbehrlich zu
sein. Gerade jetzt. Louis freute sich je länger je mehr, ihn
kennenzulernen.

Die Küche räumten sie zu dritt auf.
